% =====================================================================
%  1bat-ud5-15.tex  —  SESSIÓ 15: Taller de preparació de la prova
%  (sense preàmbul: s'inclou des de main.tex)
% =====================================================================
\horatitol{Sessió 15 — Taller de preparació de la prova}%
{Conèixer l'estructura de la prova i la rúbrica, assajar un simulacre per blocs i fer un pla d'estudi personal.}

\subsection{Idea clau}

Demà tindrem el repàs ben fresc; avui ens hi preparem \textbf{sense sorpreses}. La
prova de la sessió següent té \textbf{tres blocs}, un per cada criteri que hem
treballat (C1, C2 i C3). Saber \emph{exactament} què es demana i com es corregeix fa
que l'examen deixi de ser imprevisible. Després assajarem amb un \textbf{simulacre}.

\begin{idea}
\textbf{Estructura de la prova} (un bloc per criteri):
\begin{enumerate}[leftmargin=1.6em, itemsep=2pt]
  \item \textbf{Bloc 1 (C1) — Càlcul percentual i variacions:} factor multiplicador,
        augments i disminucions, percentatges encadenats (multiplicar índexs).
  \item \textbf{Bloc 2 (C2) — Matemàtica financera:} interès compost, TIN i TAE,
        comparar ofertes de crèdit.
  \item \textbf{Bloc 3 (C3) — Economia aplicada:} IPC i poder adquisitiu, pressupost i
        balanç.
\end{enumerate}
\textbf{Es valora}, a més del resultat: la \textbf{interpretació} dins del context
social i la \textbf{mirada crítica} davant d'ofertes i dades econòmiques.
\end{idea}

\subsection{Rúbrica simplificada}

\noindent Així es corregeix cada bloc (nivells d'assoliment):
\begin{center}
\renewcommand{\arraystretch}{1.3}
\begin{tabular}{p{2.4cm} p{10.2cm}}
\toprule
\textbf{Nivell} & \textbf{Què demostra} \\
\midrule
\textbf{AE} (excel·lent) & Calcula sense errors i \textbf{interpreta} el resultat dins
del context econòmic i social. \\
\textbf{AN} (notable) & Calcula correctament amb algun error menor; la interpretació és
essencialment bona. \\
\textbf{AS} (suficient) & Fa el procediment bàsic amb alguns errors; la interpretació
és incompleta. \\
\textbf{NA} (no assolit) & No identifica el procediment o comet errors que
n'invaliden el resultat. \\
\bottomrule
\end{tabular}
\end{center}

\subsection{Simulacre — un ítem de cada bloc}

\begin{exemple}[com es resol un ítem ben fet, de cap a cap]
\textbf{Ítem tipus (Bloc 2):} un microcrèdit té un TIN del $3\%$ mensual. Un ítem
\textbf{complet} respon: càlcul \emph{i} interpretació.
\begin{itemize}[leftmargin=1.4em, itemsep=1pt]
  \item TAE: $(1,03)^{12}-1\approx 0,426=42,6\%$.
  \item \textbf{Interpretació:} un «$3\%$» mensual és una TAE del $42,6\%$, molt cara;
        desaconsellaria el crèdit. \emph{Aquesta frase final és la que distingeix un AE
        d'un AS.}
\end{itemize}
\end{exemple}

\begin{exercicis}
\textbf{Resol el simulacre amb temps controlat. Un ítem per bloc:}
\begin{exlist}
  \item \textbf{(Bloc 1 — C1)} Un material de $\num{250}\eur$ puja un $20\%$ i, sobre el
        preu resultant, té un descompte del $10\%$. Calcula el preu final i digues a
        quina variació global equival.
  \item \textbf{(Bloc 2 — C2)} Calcula el capital final de $\num{800}\eur$ al $4\%$
        compost anual durant $3$ anys. Després, explica per què un crèdit ràpid amb un
        TIN del «$5\%$ mensual» és perillós (calcula'n la TAE).
  \item \textbf{(Bloc 3 — C3)} Un lloguer de $\num{600}\eur$ s'actualitza amb un IPC del
        $3,5\%$. Calcula la nova quota. Si l'ajut de la família ($\num{850}\eur$) es
        congela, quant li queda ara per a la resta de despeses, i quant li quedava abans?
\end{exlist}
\end{exercicis}

\subsection{Formulari-resum (suport)}

\begin{idea}
\textbf{Tingues-ho a mà mentre prepares la prova:}
\begin{itemize}[leftmargin=1.4em, itemsep=2pt]
  \item \textbf{Factor:} augment del $p\%\rightarrow 1+\frac{p}{100}$; disminució
        $\rightarrow 1-\frac{p}{100}$. Per desfer-ho, \emph{divideix}.
  \item \textbf{Encadenats:} multiplica els índexs; del producte llegeixes el $\%$
        global.
  \item \textbf{Interès compost:} $C_f=C_0\per (1+r)^t$.
  \item \textbf{TAE} (d'un TIN mensual $r$): $\text{TAE}=(1+r)^{12}-1$.
  \item \textbf{Poder adquisitiu} d'una quantitat congelada amb inflació: divideix per
        l'índex de preus.
  \item \textbf{Balanç} d'un pressupost: ingressos $-$ despeses (ha de ser positiu).
\end{itemize}
\end{idea}

\noindent\textbf{Pla d'estudi personal.} A partir de la teva graella d'autoavaluació
(sessió 5) i de com t'ha anat el simulacre, anota els \textbf{dos o tres tipus
d'exercici} que has de repassar amb més atenció abans de la prova:
\begin{center}
\renewcommand{\arraystretch}{1.5}
\begin{tabular}{c p{10.5cm}}
\toprule
\textbf{\#} & \textbf{Què he de repassar (i com)} \\
\midrule
1 & \rule{9.8cm}{0.4pt} \\
2 & \rule{9.8cm}{0.4pt} \\
3 & \rule{9.8cm}{0.4pt} \\
\bottomrule
\end{tabular}
\end{center}

% ---------------------------------------------------------------------
\fullsolucions{Sessió 15 — simulacre}
\begin{solucions}
\begin{sollist}
  \item \textbf{(Bloc 1)} $\num{250}\per 1,20\per 0,90=\num{250}\per 1,08=\num{270}\eur$.
        L'índex global és $1,08$: una \textbf{pujada global del $8\%$} (no del $10\%$ de
        restar $20-10$). El preu final és $\num{270}\eur$.
  \item \textbf{(Bloc 2)} Capital final:
        $\num{800}\per (1,04)^3\approx\num{899,89}\eur$. TAE d'un $5\%$ mensual:
        $(1,05)^{12}-1\approx 0,7959=79,6\%$. És perillós perquè el «$5\%$» sembla petit
        però és \textbf{mensual i compost}: la TAE real és gairebé un $80\%$ anual, de
        manera que el deute es dispara molt de pressa.
  \item \textbf{(Bloc 3)} Nova quota: $\num{600}\per 1,035=\num{621}\eur$. Amb l'ajut
        congelat de $\num{850}\eur$, ara li queden $\num{850}-\num{621}=\num{229}\eur$;
        abans li quedaven $\num{850}-\num{600}=\num{250}\eur$. Ha perdut $\num{21}\eur$
        al mes de poder de compra, tot i que l'ajut «no ha baixat».
\end{sollist}
\end{solucions}
